\clearpage
\ExposeTOC
\begin{center}
S G A 7

\bigskip
\underline{第 XII 讲}

\bigskip
\underline{二次超曲面}

\bigskip
\underline{P. DELIGNE}
\end{center}

\bigskip
0.0. 本讲汇集若干有关二次超曲面(quadric)的熟知结果, 我们在证明 Picard--Lefschetz 公式时将要用到它们.

0.1. 设 $m$ 为概形 $S$ 上 $\mathcal O_S$-模 $M$ 的截面. 若对每个 $s\in S$, $m$ 在由基变换 $\{s\}\longrightarrow S$ 从 $M$ 得到的 $k(s)$-向量空间 $M_s$ 中的像非零, 就称 $m$ 在 $S$ 上\underline{处处非零}(nowhere zero). 当 $S=\Spec(A)$ 而 $m$ 为 $A$-模 $M$ 的元素时, 也采用同一术语.

\medskip
1. \underline{二次型}

1.1. 设 $S=\Spec(A)$ 为仿射概形, $V$ 为秩为 $n$ 的局部自由 $A$-模, $Q$ 为 $V$ 上的二次型 (Bourbaki, \emph{Algèbre}, 第~9 章, \S3, 第~4 节). 回顾若相应的双线性型

\begin{equation*}
\Phi(x,y)=Q(x+y)-Q(x)-Q(y)
\end{equation*}

非退化, 即它诱导 $V$ 到其对偶 $V^*$ 的同构, 则称 $Q$ \underline{非退化}(nondegenerate). 若 $2$ 在 $A$ 中不可逆, 则只有当 $n$ 为偶数时才可能如此 (同上).

%%%%%%%%%% source scan idx 70 | folio 63 %%%%%%%%%%

$V$ 上的二次型构成一个投射 $A$-模; 二次型 $Q$ 在 $S$ 上``处处非零'' (0.1) 当且仅当元素 $Q(v)$ 生成 $A$ 的单位理想.

记号 $\mathbb P(V)$ 采用 EGA 的意义; 回顾若 $A$ 为域, 则 $\mathbb P(V)$ 是其有理点集为商集 $(V^*\setminus\{0\})/A^*$ 的射影空间.

二次型 $Q$ 可等同于 $S$ 上射影空间 $\mathbb P(V^*)$ 的可逆层 $\mathcal O(2)$ 的一个截面. 当 $Q$ 处处非零时, (``齐次'') 方程 $Q=0$ 定义了 $\mathbb P(V^*)$ 的一个子概形, 它在 $S$ 上平坦且纯相对维数为 $n-2$. 此概形称为\underline{由 $Q$ 定义的二次超曲面}. 若型 $Q$ 在 $S$ 上处处非零, 且其定义的二次超曲面在 $S$ 上光滑, 则称它为\underline{普通的}(ordinary).

型 $Q$ 普通当且仅当从 $A$ 到任意域作标量扩张所得的型均普通. 若 $A$ 为域, 则:

\begin{enumerate}[label=\alph*),leftmargin=2.5em]
\item 若 $n$ 为偶数或 $\operatorname{char}(A)\ne2$, 则 $Q$ 普通当且仅当 $Q$ 非退化;

\item 若 $n$ 为奇数且 $\operatorname{char}(A)=2$, 则 $Q$ 普通当且仅当相应的 (交替) 双线性型 $\Phi$ 的核 $N$ 的维数为一, 且 $Q$ 在 $N$ 上不恒为零.
\end{enumerate}

\underline{命题} 1.2. \underline{若 $Q$ 普通且 $n=2m$ (分别地, $n=2m+1$), 则在 $S$ 的平展拓扑下局部地, $V$ 有一组基 $e_1,\ldots,e_n$, 使得}

%%%%%%%%%% source scan idx 71 | folio 64 %%%%%%%%%%

\begin{equation*}
Q\!\left(\sum_{i=1}^{n}x_i e_i\right)
=\sum_{i=1}^{m}x_i x_{i+m}
\end{equation*}

\noindent\underline{(分别地,}

\begin{equation*}
Q\!\left(\sum_{i=1}^{n}x_i e_i\right)
=\sum_{i=1}^{m}x_i x_{i+m}+\lambda x_{2m+1}^{2},
\qquad (\lambda\text{ 可逆}).
\end{equation*}

\footnote{印本原文中两个求和的上限均为 $m-1$. 这会留下二维的根基, 与命题矛盾; 标准的普通型标准形以及随后的归纳均要求上限为 $m$.}

对 $m$ 归纳证明 1.2. 若 $m=0$, 断言显然成立. 若 $m>0$, 方程 $Q=0$ 的二次超曲面在 $S$ 上光滑, 且几何纤维非空; 因而在平展拓扑下局部地有截面. 这样的截面局部地由在 $\Spec(A)$ 上处处非零且满足 $Q(e)=0$ 的元素 $e\in V$ 定义. 在 $S$ 的任何点上, $e$ 都不属于 $\Phi$ 的核; 因此局部地可找到 $f'$, 使 $\Phi(e,f')=1$. 若

\begin{equation*}
f=-Q(f')e+f',
\end{equation*}

则

\begin{equation*}
Q(e)=Q(f)=0,\qquad \Phi(e,f)=1.
\end{equation*}

\footnote{印本中 $f$ 的定义的末项写作 $f$; 这是原文的明显笔误, 应为 $f'$, 所列结论也证实了这一点.}

若 $V_2$ 为 $e$ 和 $f$ 生成的子模 $V_1$ 的正交补, 则 $V=V_1\oplus V_2$, 对 $V_2$ 应用归纳假设即得结论.

1.3. 关于 $Q$ 的 Clifford 代数的定义及主要性质, 参见 Bourbaki, \emph{Algèbre}, 第~9 章, \S9. 以 $\rho$ 表示 $V$ 到 $C(Q)$ 的典范映射.

\underline{引理} 1.3.1. \underline{在 1.1 的假设下, 对每个 $\lambda\in A$, 存在唯一同态}

\begin{equation*}
[\lambda]\colon C^+(\lambda Q)\longrightarrow C^+(Q)
\end{equation*}

\underline{使得对 $x,y\in V$,}

\begin{equation*}
[\lambda]\bigl(\rho(x)\rho(y)\bigr)
=\lambda\rho(x)\rho(y).
\end{equation*}

%%%%%%%%%% source scan idx 72 | folio 65 %%%%%%%%%%

此同态由 $V$ 的张量代数的偶部 $T^+(V)$ 的自同态 $[\lambda]$ 取商诱导而得, 其中该自同态在 $T^{2r}(V)$ 上按 $\lambda^r$ 作用. 事实上,

\begin{equation*}
[\lambda]\bigl(x\otimes x-\lambda Q(x)\bigr)
=\lambda\bigl(x\otimes x-Q(x)\bigr).
\end{equation*}

假定 $Q$ 在 $S$ 上处处非零. 则引理 1.3.1 在 $\lambda\in A^*$ 对应的代数 $C^+(\lambda Q)$ 之间定义了相容的同构系.

因此, 代数 $C^+(Q)$ (在唯一同构的意义下) 只依赖于 $V$ 以及相差一个乘因子的 $Q$; 换言之, $C^+(Q)$ 只依赖于 $V$ 和 $\mathbb P(V^*)$ 中方程为 $Q=0$ 的子概形. 每个比率可逆的位似都是对 $(V,X\subset\mathbb P(V^*))$ 的自同构; 可验证这些自同构在 $C^+(V,X)$ 上的作用平凡. 因而代数 $C^+(V,X)$ 只依赖于 $A$ 上的概形对 $X\subset\mathbb P(V^*)$. 我们将其记为 $C^+(X,\mathbb P(V^*))$.
